% Part: first-order-logic
% Chapter: syntax-and-semantics
% Section: formation-sequences

\documentclass[../../../include/open-logic-section]{subfiles}

\begin{document}

\olfileid{pl}{syn}{fseq}

\olsection{Urutan Pambentukan}

Nemtokake !!{formula}s nganggo definisi induktif, lan teknik
gegandhengane kanggo mbuktekake sipat !!{formula}s kanthi induksi,
iku pendekatan kang endah lan efisien. Nanging, kita uga kadhangkala
perlu nggatekake pendekatan saka ngisor menyang ndhuwur, tahap demi
tahap, kanggo mbangun !!{formula}s. Ing kene kita nindakake iki nganggo
gagasan \emph{urutan pambentukan}.

\begin{defn}[Urutan pambentukan kanggo formula]
\ollabel{defn:fseq-frm}
Urutan winates $\tuple{!A_0,\dotsc,!A_n}$ saka rerangken simbol ing
basa~$\Lang L_0$ diarani \emph{urutan pambentukan} kanggo $!A$ yen
$!A \ident !A_n$ lan kanggo kabeh $i \leq n$, $!A_i$ iku formula
atomik utawa ana $j,k < i$ kang ndadekake salah siji saka perkara iki
lumaku:
\begin{enumerate}
    \tagitem{prvNot}{$!A_i \ident \lnot !A_j$.}{}%
    \tagitem{prvAnd}{$!A_i \ident (!A_j \land !A_k)$.}{}%
    \tagitem{prvOr}{$!A_i \ident (!A_j \lor !A_k)$.}{}%
    \tagitem{prvIf}{$!A_i \ident (!A_j \lif !A_k)$.}{}%
    \tagitem{prvIff}{$!A_i \ident (!A_j \liff !A_k)$.}{}%
\end{enumerate}
\end{defn}

\begin{ex}
\[
\tuple{
    \Obj p_0,
    \Obj p_1,
    (\Obj p_1 \land \Obj p_0),
    \lnot (\Obj p_1 \land \Obj p_0)
}
\]
iku urutan pambentukan saka
$\lnot (\Obj p_1 \land \Obj p_0)$, semono uga
\[
\tuple{
    \Obj p_0,
    \Obj p_1,
    \Obj p_0,
    (\Obj p_1 \land \Obj p_0),
    (\Obj p_0 \lif \Obj p_1),
    \lnot (\Obj p_1 \land \Obj p_0)
}.
\]
%
Kaya kang katon saka tuladha kapindho, urutan pambentukan bisa ngemot
``sampah'': formula kang mubazir utawa ora melu mbangun asil.
\end{ex}

\begin{prop}
\ollabel{prop:formed}
Saben !!{formula}~$!A$ ing~$\Frm[L_0]$ duwe urutan pambentukan.
\end{prop}

\begin{proof}
Upamakna $!A$ iku atomik. Banjur urutan~$\tuple{!A}$ minangka urutan
pambentukan kanggo~$!A$.
%
Saiki upamakna $!B$ lan~$!C$ duwe urutan pambentukan
$\tuple{!B_0,\dotsc,!B_n}$ lan $\tuple{!C_0,\dotsc,!C_m}$
miturut urutane.
%
\begin{enumerate}
  \tagitem{prvNot}{Yen $!A \ident \lnot !B$,
      mula $\tuple{!B_0,\dotsc,!B_n,\lnot !B_n}$
      iku urutan pambentukan kanggo~$!A$.}{}
  \tagitem{prvAnd}{Yen $!A \ident (!B \land !C)$,
      mula $\tuple{!B_0,\dotsc,!B_n,!C_0,\dotsc,!C_m,(!B_n \land !C_m)}$
      iku urutan pambentukan kanggo~$!A$.}{}
  \tagitem{prvOr}{Yen $!A \ident (!B \lor !C)$,
      mula $\tuple{!B_0,\dotsc,!B_n,!C_0,\dotsc,!C_m,(!B_n \lor !C_m)}$
      iku urutan pambentukan kanggo~$!A$.}{}
  \tagitem{prvIf}{Yen $!A \ident (!B \lif !C)$,
      mula $\tuple{!B_0,\dotsc,!B_n,!C_0,\dotsc,!C_m,(!B_n \lif !C_m)}$
      iku urutan pambentukan kanggo~$!A$.}{}
  \tagitem{prvIff}{Yen $!A \ident (!B \liff !C)$,
      mula $\tuple{!B_0,\dotsc,!B_n,!C_0,\dotsc,!C_m,(!B_n \liff !C_m)}$
      iku urutan pambentukan kanggo~$!A$.}{}
\end{enumerate}
Miturut asas induksi ing !!{formula}s, saben !!{formula} duwe urutan
pambentukan.
\end{proof}

Kita uga bisa mbuktekake walikane. Iki wigati amarga nuduhake yen rong
cara kita kanggo nemtokake formula iku ekuivalen: kalorone menehi asil
kang padha. Iki uga ateges kita bisa mbuktekake teorema ngenani formula
nganggo induksi lumrah ing dawa urutan pambentukan.

\begin{lem}
\ollabel{lem:fseq-init}
Upamakna $\tuple{!A_0,\dotsc,!A_n}$ iku urutan pambentukan kanggo~$!A_n$,
lan $k \leq n$. Banjur $\tuple{!A_0,\dotsc,!A_k}$ iku urutan
pambentukan kanggo~$!A_k$.
\end{lem}

\begin{proof}
Latihan.
\end{proof}

\begin{thm}
\ollabel{thm:fseq-frm-equiv}
$\Frm[L_0]$ iku himpunan kabeh rerangken simbol ing basa~$\Lang L_0$
kang duwe urutan pambentukan.
\end{thm}

\begin{proof}
$F$ dadi himpunan kabeh rerangken simbol ing basa~$\Lang L_0$ kang
duwe urutan pambentukan. Ing \olref[pl][syn][fseq]{prop:formed} kita
wis weruh yen $\Frm[L_0] \subseteq F$, mula saiki kita mbuktekake
walikane.

Upamakna $!A$ duwe urutan pambentukan $\tuple{!A_0,\dotsc,!A_n}$.
Kita mbuktekake yen $!A \in \Frm[L_0]$ kanthi induksi kuwat ing~$n$.
Hipotesis induksi kita yaiku saben rerangken simbol kang duwe urutan
pambentukan dawa $m < n$ kalebu ing~$\Frm[L_0]$. Miturut definisi
urutan pambentukan, $!A_n$ iku atomik utawa mesthi ana $j,k < n$ kang
ndadekake salah siji saka perkara iki lumaku:
\begin{enumerate}
    \tagitem{prvNot}{$!A_n \ident \lnot !A_j$.}{}%
    \tagitem{prvAnd}{$!A_n \ident (!A_j \land !A_k)$.}{}%
    \tagitem{prvOr}{$!A_n \ident (!A_j \lor !A_k)$.}{}%
    \tagitem{prvIf}{$!A_n \ident (!A_j \lif !A_k)$.}{}%
    \tagitem{prvIff}{$!A_n \ident (!A_j \liff !A_k)$.}{}%
\end{enumerate}
Saiki kita nalar miturut kasus. Yen $!A_n$ atomik, mula
$!A_n \in \Frm[L_0]$. Dene upamakna $!A \equiv
(!A_j \land !A_k)$. Miturut \olref[pl][syn][fseq]{lem:fseq-init},
$\tuple{!A_0,\dotsc,!A_j}$ lan $\tuple{!A_0,\dotsc,!A_k}$ iku
urutan pambentukan kanggo $!A_j$ lan~$!A_k$ miturut urutane. Amarga
kalorone iku suburutan wiwitan sejati saka urutan pambentukan
kanggo~$!A$, dawane kalorone kurang saka~$n$. Mula miturut hipotesis
induksi, $!A_j$ lan~$!A_k$ kalebu ing~$\Frm[L_0]$, dadi miturut
definisi !!a{formula}, $(!A_j \land !A_k)$ uga kalebu. Kasus liyane
nduweni penalaran kang padha.
\end{proof}

\end{document}
